\section{Resilience y Public Service: cómo seguir trabajando después de un incidente}

Los incidentes de seguridad y los procesos legales afectan la salud, la identidad y la carrera de empleados, investigadores, usuarios y líderes. La \term{resilience} (resiliencia) no consiste en fingir que no hubo daño, sino en reconstruir de forma gradual una capacidad de servicio sostenible mediante apoyo, recuperación, reflexión y reincorporación. El curso menciona el trabajo de Sullivan en un proyecto de servicio público en Ucrania; esta parte no puede tomarse como prueba del caso, pero aporta una perspectiva importante sobre el sistema humano.

\subsection{Human recovery también es parte del incident lifecycle}

El postmortem tradicional se centra en la root cause, el control gap y el owner, pero puede pasar por alto las guardias prolongadas, la presión pública, el daño moral y la carga familiar. Si una organización solo exige «restablecer el servicio», agota al personal clave antes del siguiente incidente. El recovery plan debe incluir rotación, apoyo psicológico, asistencia legal y de HR, licencias, ajustes de funciones y una vía para reincorporarse al trabajo.

\lecturefigure{14-resilience-loop.png}{La resiliencia surge de un ciclo de apoyo, recuperación, reconstrucción del sentido y regreso al servicio.}{Subtítulos locales 27:00--29:20; redibujo conceptual.}

\begin{knowledgebox}{Lectura de la figura: recuperarse no es volver al estado previo al incidente}
El incident/case provoca un impacto en la operación y en las personas; el support aporta family, peer, professional y legal care; la recovery requiere rest, perspective y new routines; service/learning convierte la experiencia en community work, mentoring y better systems. Las flechas no garantizan que cada persona se recupere de forma lineal, sino que recuerdan a la organización que debe dejar margen para ritmos distintos. Interpretar la resilience como «la persona debe ser más fuerte» oculta la responsabilidad institucional.
\end{knowledgebox}

\begin{warningbox}{No usar el servicio público ni el crecimiento personal como blanqueo legal}
Las buenas acciones de una persona después de un incidente pueden explicar sus decisiones de valores y su proceso de recuperación, pero no modifican por sí mismas los hechos ya juzgados; del mismo modo, la responsabilidad legal no debe llevar a la organización a descuidar la salud y la dimensión humana. Las notas del curso deben permitir que ambos juicios coexistan, en lugar de sustituir el análisis de la evidencia por un relato inspirador.
\end{warningbox}

\subsection{Government technical talent}

El sector público necesita talento técnico que entienda cloud, identity, AI, supply chain e incident operations; de lo contrario, las normas pueden describir objetivos sin poder evaluar su implementación. Que los líderes técnicos participen en el government no significa solo ocupar un cargo de tiempo completo: también incluye standards work, advisory groups, comentarios públicos, incident coordination y periodos cortos de servicio. Los participantes deben gestionar los conflictos de interés y anteponer la responsabilidad pública al lobbying de su empresa.

\lecturefigure{15-government-talent.png}{El talento técnico solo mejora la calidad de la gobernanza pública si permanece dentro del ciclo de retroalimentación entre políticas, implementación e industria.}{Subtítulos locales 36:20--39:17; redibujo conceptual.}

\begin{knowledgebox}{Lectura de la figura: la circulación del talento debe generar memoria institucional}
La technical expertise ayuda a describir systems, threats y operational cost; el policy design traduce los objetivos en rules, incentives y oversight; la implementation construye agency capability, metrics y enforcement; el industry feedback devuelve evidence, failure modes y burden. Si los expertos solo «aterrizan» por poco tiempo sin documentación, career path ni permanent staff, el conocimiento se va con cada persona; si la industria solo aporta posturas y no datos verificables, tampoco pueden formarse normas de alta calidad.
\end{knowledgebox}

\teachervoice{Al final, el ponente pidió a los profesionales técnicos que no consideren al government como un problema ajeno. Nota de clase: cuando los sistemas de software sostienen la vida pública, quienes entienden esos sistemas deben ayudar a diseñar las restricciones públicas; de lo contrario, las normas aparecerán de todos modos, solo que con mayor probabilidad en condiciones de información insuficiente y bajo la presión de una crisis.}

\subsection{Resumen de la sección}

La resilience incorpora la recuperación personal a la capacidad de seguridad a largo plazo, y el public service devuelve la experiencia técnica a la retroalimentación institucional. Ambos exigen pasar del relato heroico a los sistemas de apoyo y a la memoria organizacional.
